\section{Quantitative certificates: probabilities, values, and adversaries}
\label{sec:quantitative}

Every worker so far returns a verdict: the property holds, or here is a
counterexample. Three contributing efforts attach \emph{numbers} to unbounded
behaviour instead --- an expected running time, an optimal transport cost, the
value of a game against an adversary, the probability that a loop terminates.
The certificates stay exact and rational throughout, and where a fragment would
force an irrational value the efforts exclude the fragment rather than
approximate it.

\subsection{Why an invariant equation is not a value}

The natural first attempt fails, and it fails the same way in two of the three
efforts, which found it independently on the same one-state example. Suppose a
Markov chain has transition matrix $P$ and one asks for the probability $u$ of
eventually reaching a target. It satisfies the harmonic equation $u=Pu$ with
the obvious boundary conditions --- and so does the constant function $1$, and
so does a whole family of others. A single self-loop is enough to make the
system underdetermined. Solving $u=Pu$ and reporting the answer is not a proof
of anything.

What is missing is a \emph{progress witness}, and the two efforts supply
different ones for different reasons:

\begin{itemize}[leftmargin=1.4em]
\item \src{r7} adds a rank vector $r$ on the non-target, non-dead states
together with a policy-closed dead set, giving a bound on the expected time to
leave. The rank is auxiliary --- the effort repeatedly warns that it must not
be reported as a value --- and the conclusion is an exact supremum over
schedulers.
\item \src{r9} adds a \emph{properness potential} $h$ satisfying
$h(s)=1+\sum_t P_{st}h(t)$. This does double duty: $h$ \emph{is} the expected
absorption time, and it simultaneously licenses the reachability and cost
solutions. One certificate yields three exact values.
\end{itemize}

The second form is the more Lean-ready of the two, because it comes with an
explicit error ladder rather than an asymptotic remark. Where one effort proves
$\Prob_s(\tau>k)\le r(s)/k$ and observes that it vanishes, the other gives
$0\le v(s)-C_N\le V\,h(s)/(N+1)$ and its siblings as rational inequalities, one
per natural $N$ --- finite statements, provable by induction, with no limit in
the proposition.

\subsection{The correction: expectation is not a homomorphism}

\label{sec:moment-correction}
This section contains the only claim anywhere in the collection that a rule
already stated in this document would, if extended in the obvious way, prove
something false. It is worth the space.

Section~\ref{sec:closure}'s ideal-closure rule works because pullback along a
deterministic update is a \emph{ring homomorphism}: $p\mapsto p\circ T$
satisfies $(pq)\circ T=(p\circ T)(q\circ T)$, and that is exactly what licenses
multiplying a certificate identity through by a polynomial multiplier.
Expectation is not a ring homomorphism. $\E[XY]\ne\E[X]\,\E[Y]$ in general, and
the failure is not a corner case --- it is the entire content of correlation.

So a worker that keeps the shape of the ideal-closure rule while replacing
pathwise pullback by expectation is unsound, and would certify false
statements about products of random quantities. \src{r7}'s remedy is a type
distinction rather than a warning comment: the transfer interface splits into
a \code{PathwisePullback}, which carries a homomorphism field and may be used
multiplicatively, and a \code{MomentTransfer}, which is linear and positive and
\emph{has no homomorphism field at all}, so the multiplicative step cannot be
written. This is the same move as
Section~\ref{sec:trust}'s: make the unsound operation unrepresentable rather
than forbidden by prose.

\begin{invariant}[Moment transfer is linear, not multiplicative]
\label{inv:moment}
An expectation operator may be used to transport a linear combination of
observables and to preserve nonnegativity. It may not be used to transport a
product, and no interface should offer a field that would let it.
\end{invariant}

\subsection{Adversaries, and why a global descent measure is the wrong shape}

Two of the efforts handle games, and between them they correct the
cyclic-descent compiler of Section~\ref{sec:cycles} in two independent ways.

\src{r8} solves parity games with arbitrary priorities, returning a positional
strategy together with threshold-local ranks --- one rank per unfavourable
priority rather than one global measure. \src{r9} restricts to B\"uchi and
co-B\"uchi, which buys a simpler object: a rank that \emph{resets} on the
accepting set for the positive side, and a nonincreasing visit budget
$\kappa$ for the negative side, so that an accepted negative certificate
bounds the number of accepting visits rather than merely denying them.

The correction they share is this. A winning strategy in a game with a
recurrence condition produces plays that visit the good set infinitely often.
There is no globally strictly decreasing measure along such a play --- there
cannot be, since the play is infinite and the behaviour is correct. A compiler
that insists on well-founded descent everywhere therefore rejects a correct
proof. The repair is local descent with a permitted reset, which is precisely
what the threshold-local and resetting ranks provide.

\src{r8} makes the same point about probabilistic termination from the other
direction: a loop can terminate with probability one while having individual
executions that increase any given measure. Demanding a pointwise decrease is
demanding something false. The correct local obligation is a decrease
\emph{in expectation}, and it should be routed to the polynomial machinery of
Section~\ref{sec:nonlinear} in that form --- never as the stronger pointwise
condition.

\begin{remark}[One effort specifies what another implements]
The drift obligation $V(x)-\sum_\xi p_\xi V(F(x,\xi))\ge\varepsilon>0$ with
$V\ge0$ appears in \src{r8} as a designed composition, explicitly not executed
by its finite-state prototype. It appears in \src{r7} as an implemented worker
with a completeness theorem for the polynomial identity solve. Their worked
examples coincide numerically: the biased walk with $p=2/3$ has potential
$V(n)=3n$ in both. Neither effort cites the other, and this document is the
first place the two halves appear together.
\end{remark}

\subsection{Couplings, duals, and a finite object for an infinite comparison}

The third family is relational. To compare two probabilistic systems one
exhibits a \emph{coupling}: a joint distribution whose marginals are the two
given distributions. The certificate is the joint table, and the checker sums
rows and columns --- there is nothing subtle in the checking, which is the
point. Optimality, when claimed, is certified by a dual: a Farkas witness with
matching objective value for transport cost, or a Hall-type violated set when
no coupling exists at all.

Both efforts that build couplings also build the \emph{ordered deletion trace}
needed when a relation has to be pruned to a largest fixed point, and both are
careful about the same trap: a negative simulation result is not a trace
inequality. A failure to build a coupling at one pair says the relation does
not hold; converting that into a statement about the systems' observable
behaviour requires a separate argument, and one effort routes such a result to
the weighted-word lane of Section~\ref{sec:closure} rather than terminating the
search with a refutation.

\subsection{A finite certificate for an unbounded stack}

\src{r9} adds a model the rest of the collection does not have: a pushdown
system, where the configuration space is infinite because the stack is. Its
positive certificate is a set of \emph{pop summaries} --- what the system can do
between pushing a symbol and popping it --- assembled into a proof DAG with
shared subproofs, and its negative certificate is a closed exclusion relation.

The compression this buys is the same phenomenon Section~\ref{sec:closure}'s
multilinear lane exhibits on trees, and both are worth seeing: a DAG of
\emph{31 nodes} represents a run of $2^{31}-1$ steps. The article states plainly
that this length is \emph{represented, not executed} --- concrete expansion was
checked only through depth ten --- and that distinction is the difference
between a compression result and a fabricated benchmark.

\begin{remark}[What a finite summary does and does not settle]
\src{r9} is also the effort that states the limit of the whole approach most
precisely, and it does so against its own interests. A weighted-word identity
concerns products along \emph{finite} words; no amount of such evidence
establishes that an infinite-horizon limit has the claimed boundary values.
That is exactly why the properness witness above is not optional. And combining
the two directions is not free: a probabilistic pushdown process can produce
nonlinear probability equations and is not automatically a finite rational
Markov chain. That sentence is the boundary of this entire section.
\end{remark}

\subsection{What is excluded --- and a later correction to that exclusion}

All three efforts are exact-rational and none has a tolerance parameter. The
fragments they exclude are, in every case, the fragments where the answer would
stop being rational: stochastic games with nested minimum/maximum expectations,
continuous distributions, interval-uncertain probabilities, symbolic probability
parameters, unbounded-state chains, and probabilistic pushdown. These are not
gaps in an implementation. They are the boundary of exactness, and drawing it
explicitly is what lets everything inside it be checked by comparing rationals.

A later effort accepts that as an implementation boundary and denies that it is
the boundary of \emph{certification} \src{s2}:

\begin{quote}\small
That exclusion is a useful implementation boundary, but it is not a boundary of
\emph{exact certification}: rational enclosures can certify an irrational
value, and a finite polynomial root description can sometimes identify it
exactly.
\end{quote}

The point is well taken and the correction is adopted here. An irrational
quantity can be pinned by a certificate that is entirely rational --- a
descending sequence of rational upper bounds together with an ascending
sequence of lower bounds, each step checkable, and the target
inequality decided as soon as the interval clears it. Nothing needs to be named.

\subsection{Certifying a value that is not rational}
\status{Implemented and tested in Python \src{s2}; no Lean discharge.}

The worked case is a probabilistic recursive system: a vector of polynomial
equations $x=P(x)$ with nonnegative rational coefficients whose least fixed
point in $[0,1]^n$ is the extinction or termination probability. That least
fixed point is generally irrational, and it is generally \emph{not} the
obvious solution.

\begin{remark}[The fixed point you find need not be the one you want]
For $P(x)=\tfrac14+\tfrac34x^3$ the value $x=1$ satisfies $P(x)=x$ exactly.
The least fixed point is $(\sqrt{21}-3)/6\approx0.2637626158$. Reporting $1$
because it solves the equation would report a termination probability that is
wrong by a factor of nearly four, and no amount of checking the equation
detects it. This is the same failure as
\S\ref{sec:quantitative}'s harmonic equation, one degree up.
\end{remark}

Two certificates bracket the answer. An upper bound is a \emph{prefixed
point}: any $u$ with $P(u)\le u$ satisfies $q\le u$, by induction on the
Kleene iterates --- one polynomial evaluation and a comparison. A lower bound
needs more care, because the obvious Newton step requires inverting
$I-J(\ell)$ and rounding the result can break the inequality: $I-J$ has
negative off-diagonal entries, so componentwise rounding-down of a solution of
$(I-J)d=r$ need not preserve $(I-J)d\le r$.

The repair is to avoid the inverse entirely. Supply a positive rational weight
vector $w$ with $J(\ell)w<w$ and a step $d$ with $(I-J(\ell))d\le
P(\ell)-\ell$; then every $\ell'$ between $\ell$ and $\ell+d$ is still below
the least fixed point. The checker multiplies matrices and compares --- no
linear solve, no inversion, no rounding decision.

\begin{remark}[Strictness is not decoration]
The weight condition must be strict. For $P(x)=x$ take $\ell=0$, $w=1$,
$d=1$: the residual inequality reads $0\le0$, and a checker that accepted a
non-strict weight test would certify $\ell'=1$ as a lower bound for a least
fixed point that is zero. A stored rejection control exercises exactly this
forgery.
\end{remark}

The same effort supplies a critical-case extinction certificate --- a positive
rational weight proving extinction probability one --- with its own negative
control showing why a spectral slogan is not enough: for $P(x)=x$ the matrix is
$(1)$, every positive $w$ satisfies $Mw=w$, the dependency graph is strongly
connected, and the least fixed point is zero. Without the extra nondegeneracy
condition a plausible critical-extinction checker proves something false.

\begin{remark}[Probability one is not finite expected work]
An extinction certificate concludes that termination happens with probability
one. It does not bound the expected number of steps, and the effort names its
receipt accordingly rather than calling it \code{terminates}. That is the same
distinction Section~\ref{sec:wsts} draws between a terminating algorithm and a
bounded one.
\end{remark}

Two further limits deserve the same visibility. Rational bit complexity is
additional to every operation count reported anywhere in this document ---
denominators grow, and a small matrix dimension does not imply small numbers.
And a certificate that bounds graph steps bounds graph steps: a source adapter
whose steps have variable duration needs its own timing argument before any of
these numbers becomes a statement about time.
